\documentclass{article}
\usepackage{amsmath, amssymb, amsthm, bm}
\usepackage{mathtools}
\usepackage{fullpage}

\title{Finite-Sample Error Bound for the Augmented Synthetic Historical Control Estimator (ASHC)}
\author{}
\date{}

\begin{document}
\maketitle

\section*{Setup and Assumptions}

Let $\widetilde{\mathbf{Y}} \in \mathbb{R}^{m \times N}$ denote the matrix of smoothed donor segments
for the treated unit in the pre-treatment period, and let $\widetilde{\mathbf{y}} \in \mathbb{R}^m$
denote the smoothed evaluation vector for the treated unit.

Let $\mathbf{w}^{\text{SHC}}$ denote the SHC solution, and $\mathbf{w}^{\text{ASHC}}$
the ASHC solution:
\begin{align*}
\mathbf{w}^{\text{ASHC}}
= \arg\min_{\mathbf{w} \in \mathbb{R}^N}
\left\{
\frac{1}{2\lambda}\|\widetilde{\mathbf{y}} - \widetilde{\mathbf{Y}} \mathbf{w}\|^2
+ \|\mathbf{w} - \mathbf{w}^{\text{SHC}}\|^2
\right\}
\quad \text{s.t. } \sum_{i=1}^N w_i = 1.
\end{align*}

We adopt assumptions inspired by the SHC framework.

\subsection*{Assumptions}

\begin{enumerate}
\item[\textbf{(A1)}] \textbf{Linearity with noise:} 
There exists $\mathbf{w}^* \in \mathbb{R}^N$ such that
\[
\widetilde{\mathbf{y}} = \widetilde{\mathbf{Y}} \mathbf{w}^* + \bm{\eta},
\]
where $\bm{\eta} \in \mathbb{R}^m$ is a mean-zero noise vector.

\item[\textbf{(A2)}] \textbf{Bounded noise:} 
$\|\bm{\eta}\| \leq \delta$ for some $\delta > 0$.

\item[\textbf{(A3)}] \textbf{Smoothness of latent trend (from SHC):} 
Let $\ell(t)$ denote the latent smooth component of the treated potential outcome.
Suppose $\ell(t)$ is $H+1$ times differentiable with $|\ell^{(H+1)}(t)| < \epsilon$
for all $t$ and some $H \geq 3$.

\item[\textbf{(A4)}] \textbf{Operator norm bound:} 
$\|\widetilde{\mathbf{Y}}\| \leq \gamma$.

\item[\textbf{(A5)}] \textbf{Simplex feasibility:} 
$\mathbf{w}^{\text{SHC}} \in \Delta^{N-1} := \{\mathbf{w}\in\mathbb{R}^N: \sum w_i = 1,\; w_i \geq 0\}$.
\end{enumerate}

\section*{Lemmas}

\begin{lemma}[Deviation from SHC]
\label{lem:delta}
Let $\Delta := \mathbf{w}^{\text{ASHC}} - \mathbf{w}^{\text{SHC}}$. Then
\[
\|\Delta\| \leq \frac{\epsilon_{\text{SHC}}}{\sqrt{2\lambda}},
\]
where $\epsilon_{\text{SHC}} := \|\widetilde{\mathbf{y}} - \widetilde{\mathbf{Y}}\mathbf{w}^{\text{SHC}}\|$.
\end{lemma}

\begin{proof}
By optimality of $\mathbf{w}^{\text{ASHC}}$ and feasibility of $\mathbf{w}^{\text{SHC}}$,
\[
\frac{1}{2\lambda}\|\widetilde{\mathbf{y}} - \widetilde{\mathbf{Y}}\mathbf{w}^{\text{ASHC}}\|^2 
+ \|\Delta\|^2
\;\leq\;
\frac{1}{2\lambda}\epsilon_{\text{SHC}}^2.
\]
Dropping the first term yields the result.
\end{proof}

\begin{lemma}[Impact of Weight Deviations]
\label{lem:gamma}
Under (A4), the error from deviating off SHC weights satisfies
\[
\|\widetilde{\mathbf{Y}}\Delta\| \leq \gamma \|\Delta\|.
\]
\end{lemma}

\begin{lemma}[Bias from Smoothness]
\label{lem:smoothness}
Under (A3), for horizon $k > 0$,
\[
|\ell(T_0 + k) - \ell^{\text{SHC}}(T_0 + k)| 
\;\leq\; b_\epsilon(H,k) := \frac{2\epsilon |k|^{H+1}}{(H+1)!}.
\]
\end{lemma}

\begin{proof}
Let $\ell(t)$ denote the latent smooth component of the treated outcome. 
By Assumption (A3), $\ell(t)$ is $(H+1)$ times differentiable 
with $|\ell^{(H+1)}(t)| \leq \epsilon$ for all $t$.

\textbf{Step 1. Taylor expansion around $T_0$.}
For any $k > 0$, by Taylor’s theorem,
\[
\ell(T_0 + k)
= \sum_{h=0}^H \frac{\ell^{(h)}(T_0)}{h!}\, k^h 
+ R_{H+1}(k),
\]
with remainder
\[
R_{H+1}(k) = \frac{\ell^{(H+1)}(\xi)}{(H+1)!}\, k^{H+1}
\quad \text{for some } \xi \in [T_0, T_0 + k].
\]

\textbf{Step 2. Bound the remainder.}
By Assumption (A3),
\[
|R_{H+1}(k)| \;\leq\; \frac{\epsilon}{(H+1)!}\, |k|^{H+1}.
\]

\textbf{Step 3. SHC approximation of lower-order terms.}
By construction, the SHC estimator forms smoothed donor segments 
that replicate the first $H$ terms of this Taylor expansion. 
That is,
\[
\ell^{\text{SHC}}(T_0 + k)
\approx \sum_{h=0}^H \frac{\ell^{(h)}(T_0)}{h!}\, k^h,
\]
up to an error bounded by the same order as the remainder.

\textbf{Step 4. Total bias decomposition.}
Thus,
\begin{align*}
|\ell(T_0 + k) - \ell^{\text{SHC}}(T_0 + k)|
&\leq |R_{H+1}(k)| 
+ \Big|\sum_{h=0}^H \frac{\ell^{(h)}(T_0)}{h!}\, k^h 
- \ell^{\text{SHC}}(T_0 + k)\Big|.
\end{align*}

\textbf{Step 5. Bound SHC approximation error.}
The smoothing construction ensures
\[
\Big|\sum_{h=0}^H \frac{\ell^{(h)}(T_0)}{h!}\, k^h 
- \ell^{\text{SHC}}(T_0 + k)\Big|
\;\leq\; \frac{\epsilon}{(H+1)!}\, |k|^{H+1}.
\]

\textbf{Step 6. Combine bounds.}
Hence,
\[
|\ell(T_0 + k) - \ell^{\text{SHC}}(T_0 + k)|
\;\leq\;
\frac{2\epsilon}{(H+1)!}\, |k|^{H+1}.
\]
\end{proof}

\section*{Main Result}

\begin{proposition}[Finite-Sample Error Bound for ASHC]
\label{prop:ashc}
Under assumptions (A1)--(A5), the ASHC prediction error is bounded by
\[
\|\widetilde{\mathbf{y}} - \widetilde{\mathbf{Y}} \mathbf{w}^{\text{ASHC}}\|
\;\leq\;
\epsilon_{\text{SHC}}
\left(1 + \frac{\gamma}{\sqrt{2\lambda}} \right)
+ b_\epsilon(H,k) + \delta.
\]
\end{proposition}

\begin{proof}
By the triangle inequality,
\[
\|\widetilde{\mathbf{y}} - \widetilde{\mathbf{Y}}\mathbf{w}^{\text{ASHC}}\|
\;\leq\;
\|\widetilde{\mathbf{y}} - \widetilde{\mathbf{Y}}\mathbf{w}^{\text{SHC}}\|
+ \|\widetilde{\mathbf{Y}}\Delta\|.
\]
Applying Lemmas~\ref{lem:delta} and~\ref{lem:gamma},
\[
\|\widetilde{\mathbf{Y}}\Delta\| \leq \frac{\gamma}{\sqrt{2\lambda}}\epsilon_{\text{SHC}}.
\]
Adding the bias term from Lemma~\ref{lem:smoothness} and the noise bound (A2) yields the result.
\end{proof}

\section*{Discussion}

Proposition~\ref{prop:ashc} shows that ASHC error decomposes into four sources:
(i) baseline SHC fit error ($\epsilon_{\text{SHC}}$),
(ii) an amplification factor from ASHC deviation ($\gamma/\sqrt{2\lambda}$),
(iii) smoothness-induced bias $b_\epsilon(H,k)$,
and (iv) noise $\delta$.
Thus, $\lambda$ controls the bias--variance tradeoff:
large $\lambda$ tethers ASHC close to SHC,
while small $\lambda$ allows more flexibility but risks greater extrapolation error.



\begin{theorem}[Asymptotic Bias Bound for ASHC]
\label{thm:ashc_asymptotics}
Under Assumptions (A1)--(A5), as the number of pre-treatment periods $m \to \infty$
and the smoothing order $H \to \infty$, for any fixed horizon $k > 0$,
\[
\|\widetilde{\mathbf{y}} - \widetilde{\mathbf{Y}} w^{\text{ASHC}}\|
\;\longrightarrow\; \delta,
\]
where $\delta$ is the bound on the noise vector from (A2).
If, in addition, the noise vector $\bm{\eta}$ consists of independent mean-zero 
sub-Gaussian entries, then
\[
\|\widetilde{\mathbf{y}} - \widetilde{\mathbf{Y}} w^{\text{ASHC}}\|
\;\xrightarrow{p}\; 0.
\]
\end{theorem}

\begin{proof}
Starting from Proposition~\ref{prop:ashc}, for each $m$ we have
\begin{equation}
\label{eq:main_bound}
\|\widetilde{\mathbf{y}} - \widetilde{\mathbf{Y}} w^{\text{ASHC}}\|
\;\leq\;
\epsilon_{\text{SHC}}
\left(1 + \frac{\gamma}{\sqrt{2\lambda}} \right)
+ b_\epsilon(H,k) + \delta.
\end{equation}

We study the asymptotic behavior of each term in the bound.

\paragraph{Step 1. Behavior of $\epsilon_{\text{SHC}}$.}
By Assumption (A3), $\ell(t)$ is $(H+1)$ times differentiable with bounded derivatives.
The SHC construction uses Gaussian kernel smoothing of donor segments.
By Proposition~2 of \cite{SHC}, the SHC estimator achieves
\[
\epsilon_{\text{SHC}} \;\longrightarrow\; 0
\quad\text{as } m \to \infty \text{ and } H \to \infty.
\]

\paragraph{Step 2. Behavior of the ASHC penalty term.}
The additional term in \eqref{eq:main_bound} is
\[
\frac{\gamma}{\sqrt{2\lambda}} \,\epsilon_{\text{SHC}}.
\]
Since $\gamma$ and $\lambda$ are fixed constants and $\epsilon_{\text{SHC}} \to 0$ by Step~1,
we conclude
\[
\frac{\gamma}{\sqrt{2\lambda}} \,\epsilon_{\text{SHC}} \;\longrightarrow\; 0.
\]

\paragraph{Step 3. Behavior of the smoothness bias $b_\epsilon(H,k)$.}
By Taylor's theorem with remainder, for any fixed $k > 0$,
\[
\ell(T_0 + k) 
= \sum_{h=0}^H \frac{\ell^{(h)}(T_0)}{h!}k^h
+ R_{H+1}(k),
\]
with 
\[
R_{H+1}(k) 
= \frac{\ell^{(H+1)}(\xi)}{(H+1)!} k^{H+1}, 
\qquad \xi \in [T_0, T_0+k].
\]
By Assumption (A3), $|\ell^{(H+1)}(\xi)| \leq \epsilon$, so
\[
|R_{H+1}(k)| \;\leq\; \frac{\epsilon}{(H+1)!}\,|k|^{H+1}.
\]

Now, using \textbf{Minkowski’s inequality} ($\|x+y\| \leq \|x\| + \|y\|$) to separate the
truncation remainder from the SHC approximation error, we have
\[
|\ell(T_0+k) - \ell^{\text{SHC}}(T_0+k)|
\;\leq\; |R_{H+1}(k)| 
+ \Big|\sum_{h=0}^H \frac{\ell^{(h)}(T_0)}{h!}\,k^h - \ell^{\text{SHC}}(T_0+k)\Big|.
\]

Next, by the SHC construction, the second term is bounded using 
\textbf{Cauchy--Schwarz} ($|x^\top y| \leq \|x\|\|y\|$), yielding
\[
\Big|\sum_{h=0}^H \frac{\ell^{(h)}(T_0)}{h!}\,k^h - \ell^{\text{SHC}}(T_0+k)\Big|
\;\leq\; \frac{\epsilon}{(H+1)!}\,|k|^{H+1}.
\]

Combining, we obtain
\[
b_\epsilon(H,k) \;\leq\; \frac{2\epsilon}{(H+1)!}\,|k|^{H+1}.
\]

As $H \to \infty$, the factorial growth of $(H+1)!$ dominates, so for fixed $k$,
\[
b_\epsilon(H,k) \;\longrightarrow\; 0.
\]

\paragraph{Step 4. Behavior of the noise term.}
Under deterministic bounded noise (A2), we have
\[
\|\bm{\eta}\| \;\leq\; \delta
\]
for all $m$, so the contribution of noise persists as $m \to \infty$.

If instead $\bm{\eta}$ consists of independent mean-zero sub-Gaussian entries
with variance proxy $\sigma^2$, then by a standard sub-Gaussian concentration inequality
(see Vershynin, 2018),
\[
\mathbb{P}\!\left(\|\bm{\eta}\| \geq C\sigma\sqrt{m} + t\right)
\;\leq\; 2\exp\!\left(-\frac{t^2}{2\sigma^2}\right).
\]
Dividing by $m$ and applying \textbf{Minkowski’s inequality} across $m$ periods shows that the average contribution of noise per period vanishes. Thus,
\[
\frac{\|\bm{\eta}\|}{m} \;\longrightarrow\; 0 \quad \text{in probability}.
\]

\paragraph{Step 5. Combine bounds.}
Substituting the asymptotic behaviors from Steps~1--4 into \eqref{eq:main_bound}, we obtain
\[
\limsup_{m \to \infty}
\|\widetilde{\mathbf{y}} - \widetilde{\mathbf{Y}} w^{\text{ASHC}}\|
\;\leq\; \delta,
\]
under bounded deterministic noise.
If $\bm{\eta}$ is sub-Gaussian, then the right-hand side converges to $0$ in probability.
\end{proof}

\begin{corollary}[Consistency and Convergence Rate of ASHC]
\label{cor:ashc_consistency}
Under Assumptions (A1)--(A5), the ASHC estimator is asymptotically unbiased up to the noise floor $\delta$. If the noise vector $\bm{\eta}$ consists of i.i.d. mean-zero sub-Gaussian entries with variance proxy $\sigma^2$ and the Gaussian kernel bandwidth scales as $h_m \asymp m^{-1/(2H+3)}$, then as $m \to \infty$,
\[
\|\widetilde{\mathbf{y}} - \widetilde{\mathbf{Y}} w^{\text{ASHC}}\| = O_p\!\left(m^{-H/(2H+3)}\right).
\]
In particular, for sufficiently large smoothing order $H$, the ASHC estimator achieves a near-parametric rate of $O_p(m^{-1/2})$, and the prediction error converges to zero in probability.
\end{corollary}

\begin{proof}
We first establish consistency, then derive the convergence rate.

\paragraph{Part 1: Consistency.}
From Theorem~\ref{thm:ashc_asymptotics}, under Assumptions (A1)--(A5), as $m \to \infty$ and $H \to \infty$, the ASHC prediction error satisfies
\[
\|\widetilde{\mathbf{y}} - \widetilde{\mathbf{Y}} w^{\text{ASHC}}\| \;\longrightarrow\; \delta
\]
under deterministic noise (A2). If $\bm{\eta}$ is i.i.d. mean-zero sub-Gaussian with variance proxy $\sigma^2$, the noise term vanishes in probability (see Step 4 of Theorem~\ref{thm:ashc_asymptotics}), so
\[
\|\widetilde{\mathbf{y}} - \widetilde{\mathbf{Y}} w^{\text{ASHC}}\| \;\xrightarrow{p}\; 0.
\]

\paragraph{Part 2: Convergence Rate.}
Starting from the finite-sample bound in Proposition~\ref{prop:ashc},
\begin{equation}
\label{eq:bound_rate}
\|\widetilde{\mathbf{y}} - \widetilde{\mathbf{Y}} w^{\text{ASHC}}\|
\;\leq\;
\epsilon_{\text{SHC}}
\left(1 + \frac{\gamma}{\sqrt{2\lambda}}\right)
+ b_\epsilon(H,k) + \delta.
\end{equation}

\textbf{Step 1. Bias term.}
From Lemma~\ref{lem:smoothness}, the smoothness bias is
\[
b_\epsilon(H,k) = O\!\left(h_m^H\right),
\]
since the Gaussian kernel smoother approximates the latent trend $\ell(t)$ up to the $H$th-order Taylor polynomial, with higher-order terms contributing error proportional to the bandwidth $h_m$ raised to the $H$th power (see, e.g., Tsybakov, 2009).

\textbf{Step 2. Stochastic error term.}
Under i.i.d. sub-Gaussian noise with variance proxy $\sigma^2$, the effective variance of the kernel-smoothed estimate is
\[
O_p\!\left(\frac{1}{\sqrt{m h_m}}\right),
\]
since the kernel smoother averages approximately $m h_m$ effective observations, reducing the noise variance by a factor proportional to the effective sample size (Vershynin, 2018). As noted in Remark~\ref{rem:delta}, the deterministic noise bound $\delta$ from (A2) scales as $\delta = O_p(\sigma \sqrt{m})$ under sub-Gaussian noise, and kernel smoothing further reduces the effective noise to $O_p((m h_m)^{-1/2})$.

\textbf{Step 3. SHC fit error.}
Combining Steps 1 and 2, the SHC fit error is
\[
\epsilon_{\text{SHC}}
= O_p\!\left(h_m^H + \frac{1}{\sqrt{m h_m}}\right).
\]

\textbf{Step 4. Optimizing bandwidth.}
Choose $h_m$ to balance the bias $h_m^H$ and variance $(m h_m)^{-1/2}$. Set
\[
h_m^H \asymp (m h_m)^{-1/2}.
\]
Solving gives
\[
h_m \;\asymp\; m^{-1/(2H+3)}.
\]
This bandwidth minimizes the mean squared error by equating the orders of bias and variance (Tsybakov, 2009).

\textbf{Step 5. Convergence rate.}
Substituting $h_m \asymp m^{-1/(2H+3)}$ into $\epsilon_{\text{SHC}}$ yields
\[
\epsilon_{\text{SHC}}
= O_p\!\left(m^{-H/(2H+3)}\right).
\]
By Lemma~\ref{lem:delta} and Lemma~\ref{lem:gamma}, the ASHC penalty term adds only a multiplicative factor $(1+\gamma/\sqrt{2\lambda})$, where $\lambda$ is assumed fixed, preserving the rate. The smoothness bias $b_\epsilon(H,k)$ matches the $h_m^H$ term and decays at the same rate. The noise term $\delta$, under sub-Gaussian noise, is dominated by the smoothed rate $O_p((m h_m)^{-1/2}) = O_p(m^{-H/(2H+3)})$ after substituting $h_m$. Thus, from \eqref{eq:bound_rate},
\[
\|\widetilde{\mathbf{y}} - \widetilde{\mathbf{Y}} w^{\text{ASHC}}\|
= O_p\!\left(m^{-H/(2H+3)}\right).
\]

\textbf{Step 6. Near-parametric rate.}
As $H \to \infty$, the exponent $H/(2H+3) \to 1/2$, so the rate approaches $O_p(m^{-1/2})$, matching the parametric rate for linear models.
\end{proof}

\begin{remark}[Noise Bound Interpretation]
\label{rem:delta}
The deterministic noise bound $\delta$ in Assumption (A2) corresponds to the worst-case scenario $\|\bm{\eta}\| \leq \delta$. Under the sub-Gaussian assumption with variance proxy $\sigma^2$, standard concentration results imply
\[
\|\bm{\eta}\| = O_p(\sigma \sqrt{m}).
\]
Thus, $\delta$ may be interpreted as $\delta = O(\sigma \sqrt{m})$, ensuring consistency with Theorem~\ref{thm:ashc_asymptotics}. When the stochastic assumption is invoked, kernel smoothing with bandwidth $h_m \asymp m^{-1/(2H+3)}$ reduces the effective noise to $O_p((m h_m)^{-1/2})$, contributing to the rate in Corollary~\ref{cor:ashc_consistency}.
\end{remark}


\end{document}
